% Ngân hàng bài dummy (tự soạn để mở rộng ngân hàng): Hàm số y = ax^2 (a khác 0) (lớp 9), phần mở rộng.
% Quy ước:
%   - Mỗi bài: \begin{problem}[id] ... \end{problem}; lời giải lồng bên trong.
%   - id: duy nhất, không dấu, không đổi, không dùng lại. Dạng <lớp><phân môn>-<chủ đề>-<số>.
%   - Mỗi bước lời giải là 1 đoạn, các bước cách nhau 1 dòng trống; mỗi bước có lý do.
%   - Dòng cuối lời giải: Đáp số: $\boxed{...}$
%   - Cách giải khác: thêm \begin{solution}[cách 2].
%   - Metadata bằng comment ngay sau \begin{problem}: % nguon:, % lop:, % kc_du_kien:, % do_kho_du_kien:
%   - Đồ thị được mô tả bằng lời, không vẽ hình.
\documentclass[12pt,a4paper]{article}
\usepackage[utf8]{vietnam}
\usepackage{amsmath,amssymb}
\newenvironment{problem}[1][]{\par\medskip\noindent\textbf{Bài #1.}\ }{\par}
\newenvironment{solution}[1][]{\par\noindent\textit{Lời giải #1.}\par}{\par}

\begin{document}

% ---------------------------------------------------------------- giá trị hàm số, điểm thuộc đồ thị
\begin{problem}[9DS-HSB2-032]
% nguon: dummy (tự soạn để mở rộng ngân hàng)
% lop: 9
% kc_du_kien: 9DS-HSB2-01
% do_kho_du_kien: 1
Cho hàm số $y = f(x) = -3x^2$.

a) Tính $f(-1)$, $f(2)$ và $f\left(\dfrac{1}{3}\right)$.

b) Các điểm $A(-2; -12)$ và $B(1; 3)$ có thuộc đồ thị của hàm số hay không?

\begin{solution}
a) Thay $x = -1$ vào công thức của hàm số: $f(-1) = -3 \cdot (-1)^2 = -3 \cdot 1 = -3$.

a) Thay $x = 2$: $f(2) = -3 \cdot 2^2 = -3 \cdot 4 = -12$.

a) Thay $x = \dfrac{1}{3}$: $f\left(\dfrac{1}{3}\right) = -3 \cdot \left(\dfrac{1}{3}\right)^2 = -3 \cdot \dfrac{1}{9} = -\dfrac{1}{3}$.

b) Điểm $(x_0; y_0)$ thuộc đồ thị khi và chỉ khi $y_0 = f(x_0)$. Với $A(-2; -12)$: $f(-2) = -3 \cdot 4 = -12$, bằng tung độ của $A$, nên $A$ thuộc đồ thị.

b) Với $B(1; 3)$: $f(1) = -3 \cdot 1 = -3 \neq 3$, nên $B$ không thuộc đồ thị.

Đáp số: a) $\boxed{f(-1) = -3,\ f(2) = -12,\ f\left(\dfrac{1}{3}\right) = -\dfrac{1}{3}}$; b) $\boxed{A \text{ thuộc đồ thị; } B \text{ không thuộc đồ thị}}$
\end{solution}
\end{problem}

% ---------------------------------------------------------------- giá trị lớn nhất, nhỏ nhất; so sánh hai giá trị
\begin{problem}[9DS-HSB2-033]
% nguon: dummy (tự soạn để mở rộng ngân hàng)
% lop: 9
% kc_du_kien: 9DS-HSB2-02
% do_kho_du_kien: 1
Cho hai hàm số $y = 4x^2$ và $y = -\dfrac{3}{4}x^2$.

a) Với mỗi hàm số, giá trị của $y$ có giá trị lớn nhất hay nhỏ nhất? Giá trị đó bằng bao nhiêu và đạt được khi $x$ bằng bao nhiêu?

b) Với hàm số $y = f(x) = 4x^2$, không tính giá trị, hãy so sánh $f(-3)$ và $f(-2)$.

\begin{solution}
a) Với hàm số $y = 4x^2$: vì $x^2 \geq 0$ với mọi $x$ nên $y = 4x^2 \geq 0$, dấu bằng xảy ra khi $x = 0$. Vậy $y$ có giá trị nhỏ nhất bằng $0$ khi $x = 0$ (không có giá trị lớn nhất).

a) Với hàm số $y = -\dfrac{3}{4}x^2$: vì $x^2 \geq 0$ nên $y = -\dfrac{3}{4}x^2 \leq 0$, dấu bằng xảy ra khi $x = 0$. Vậy $y$ có giá trị lớn nhất bằng $0$ khi $x = 0$ (không có giá trị nhỏ nhất). Điều này phù hợp với việc đồ thị là parabol có đỉnh $O$ và nằm phía trên trục hoành khi $a > 0$, phía dưới trục hoành khi $a < 0$.

b) Hàm số $y = 4x^2$ có $a = 4 > 0$ nên nghịch biến khi $x < 0$.

b) Vì $-3 < -2 < 0$ nên $f(-3) > f(-2)$ (kiểm tra: $f(-3) = 36 > 16 = f(-2)$).

Đáp số: a) $\boxed{y = 4x^2 \text{ có giá trị nhỏ nhất bằng } 0 \text{ khi } x = 0;\ y = -\dfrac{3}{4}x^2 \text{ có giá trị lớn nhất bằng } 0 \text{ khi } x = 0}$; b) $\boxed{f(-3) > f(-2)}$
\end{solution}
\end{problem}

% ---------------------------------------------------------------- tìm a, tìm điểm thuộc parabol
\begin{problem}[9DS-HSB2-034]
% nguon: dummy (tự soạn để mở rộng ngân hàng)
% lop: 9
% kc_du_kien: 9DS-HSB2-01, 9DS-PT2-01
% do_kho_du_kien: 2
Cho hàm số $y = ax^2$ ($a \neq 0$) có đồ thị là parabol $(P)$ đi qua điểm $A\left(\sqrt{3}; -6\right)$.

a) Tìm $a$.

b) Tìm các điểm thuộc $(P)$ có tung độ bằng $-8$.

c) Tìm điểm thuộc $(P)$, khác gốc toạ độ $O$, có hoành độ bằng tung độ.

\begin{solution}
a) $A\left(\sqrt{3}; -6\right)$ thuộc $(P)$ nên $-6 = a \cdot \left(\sqrt{3}\right)^2 = 3a$, suy ra $a = -2$ (thoả mãn $a \neq 0$); $(P)\colon y = -2x^2$.

b) Điểm thuộc $(P)$ có tung độ $-8$ khi hoành độ $x$ thoả mãn $-2x^2 = -8$, tức là $x^2 = 4$.

b) Hai số có bình phương bằng $4$ là $2$ và $-2$, nên có hai điểm $(2; -8)$ và $(-2; -8)$.

c) Điểm có hoành độ bằng tung độ có dạng $M(x; x)$. $M$ thuộc $(P)$ khi và chỉ khi $x = -2x^2$, tức là $2x^2 + x = 0$.

c) Phân tích thành nhân tử: $x(2x + 1) = 0$, nên $x = 0$ hoặc $x = -\dfrac{1}{2}$. Loại $x = 0$ vì $M$ khác $O$; với $x = -\dfrac{1}{2}$ ta có $M\left(-\dfrac{1}{2}; -\dfrac{1}{2}\right)$.

c) Thử lại: $-2 \cdot \left(-\dfrac{1}{2}\right)^2 = -2 \cdot \dfrac{1}{4} = -\dfrac{1}{2}$ (đúng).

Đáp số: a) $\boxed{a = -2}$; b) $\boxed{(2; -8) \text{ và } (-2; -8)}$; c) $\boxed{M\left(-\dfrac{1}{2}; -\dfrac{1}{2}\right)}$
\end{solution}
\end{problem}

% ---------------------------------------------------------------- bài toán thực tế: quãng đường phanh
\begin{problem}[9DS-HSB2-035]
% nguon: dummy (tự soạn để mở rộng ngân hàng)
% lop: 9
% kc_du_kien: 9DS-HSB2-01, 9DS-PT2-01
% do_kho_du_kien: 2
Khi một ô tô đang chạy với vận tốc $v$ (m/s) thì hãm phanh, quãng đường $s$ (m) xe còn đi thêm được cho đến khi dừng hẳn được tính gần đúng bởi công thức $s = \dfrac{1}{8}v^2$.

a) Xe đang chạy với vận tốc $12$ m/s thì hãm phanh. Xe còn đi thêm bao nhiêu mét?

b) Một xe hãm phanh và đi thêm $32$ m thì dừng hẳn. Tìm vận tốc của xe lúc bắt đầu hãm phanh.

c) Khi vận tốc lúc bắt đầu hãm phanh tăng gấp đôi thì quãng đường xe còn đi thêm tăng lên mấy lần?

\begin{solution}
a) Thay $v = 12$ vào công thức: $s = \dfrac{1}{8} \cdot 12^2 = \dfrac{144}{8} = 18$ (m).

b) Cho $s = 32$: $\dfrac{1}{8}v^2 = 32$, nhân hai vế với $8$ được $v^2 = 256$.

b) Hai số có bình phương bằng $256$ là $16$ và $-16$; vì vận tốc là số dương nên $v = 16$ (m/s).

c) Gọi $v$ là vận tốc ban đầu, ứng với quãng đường $s = \dfrac{1}{8}v^2$. Khi vận tốc là $2v$ thì quãng đường là $s' = \dfrac{1}{8}(2v)^2 = \dfrac{4v^2}{8} = 4 \cdot \dfrac{1}{8}v^2 = 4s$.

c) Vậy khi vận tốc tăng gấp đôi thì quãng đường xe còn đi thêm tăng lên $4$ lần.

Đáp số: a) $\boxed{18 \text{ m}}$; b) $\boxed{16 \text{ m/s}}$; c) $\boxed{4 \text{ lần}}$
\end{solution}
\end{problem}

% ---------------------------------------------------------------- tìm a, so sánh, điểm có tung độ cho trước
\begin{problem}[9DS-HSB2-036]
% nguon: dummy (tự soạn để mở rộng ngân hàng)
% lop: 9
% kc_du_kien: 9DS-HSB2-01, 9DS-HSB2-02
% do_kho_du_kien: 2
Biết đồ thị của hàm số $y = f(x) = ax^2$ ($a \neq 0$) đi qua điểm $A(4; -8)$.

a) Tìm $a$.

b) Không tính giá trị, hãy so sánh $f(-5)$ và $f(-4)$.

c) Tìm các điểm thuộc đồ thị có tung độ bằng $-18$.

\begin{solution}
a) Đồ thị đi qua $A(4; -8)$ nên $-8 = a \cdot 4^2 = 16a$, suy ra $a = -\dfrac{1}{2}$ (thoả mãn $a \neq 0$); hàm số là $y = -\dfrac{1}{2}x^2$.

b) Vì $a = -\dfrac{1}{2} < 0$ nên hàm số đồng biến khi $x < 0$.

b) Ta có $-5 < -4 < 0$, nên $f(-5) < f(-4)$ (kiểm tra: $f(-5) = -\dfrac{25}{2}$, $f(-4) = -8$ và $-\dfrac{25}{2} < -8$).

c) Điểm có tung độ $-18$ thuộc đồ thị khi $-\dfrac{1}{2}x^2 = -18$, tức là $x^2 = 36$, suy ra $x = 6$ hoặc $x = -6$.

c) Vậy có hai điểm $(6; -18)$ và $(-6; -18)$, đối xứng nhau qua trục $Oy$.

Đáp số: a) $\boxed{a = -\dfrac{1}{2}}$; b) $\boxed{f(-5) < f(-4)}$; c) $\boxed{(6; -18) \text{ và } (-6; -18)}$
\end{solution}
\end{problem}

% ---------------------------------------------------------------- giao điểm của parabol và đường thẳng, độ dài đoạn thẳng
\begin{problem}[9DS-HSB2-037]
% nguon: dummy (tự soạn để mở rộng ngân hàng)
% lop: 9
% kc_du_kien: 9DS-HSB2-03, 9DS-PT2-02
% do_kho_du_kien: 2
Cho parabol $(P)\colon y = -\dfrac{1}{4}x^2$ và đường thẳng $(d)\colon y = x - 3$.

a) Tìm toạ độ các giao điểm $A$, $B$ của $(P)$ và $(d)$ (hoành độ của $A$ lớn hơn hoành độ của $B$).

b) Tính độ dài đoạn $AB$.

\begin{solution}
a) Hoành độ giao điểm của $(P)$ và $(d)$ là nghiệm của phương trình $-\dfrac{1}{4}x^2 = x - 3$; nhân hai vế với $-4$ rồi chuyển vế: $x^2 + 4x - 12 = 0$.

a) Dùng công thức nghiệm thu gọn với $b' = 2$: $\Delta' = 2^2 - 1 \cdot (-12) = 16$, $\sqrt{\Delta'} = 4$, nên $x_1 = -2 + 4 = 2$ và $x_2 = -2 - 4 = -6$.

a) Tính tung độ từ $y = x - 3$: với $x = 2$ thì $y = -1$; với $x = -6$ thì $y = -9$. Vậy $A(2; -1)$ và $B(-6; -9)$.

a) Kiểm tra trên $(P)$: $-\dfrac{1}{4} \cdot 2^2 = -1$ và $-\dfrac{1}{4} \cdot (-6)^2 = -9$ (đúng).

b) Theo công thức khoảng cách giữa hai điểm: $AB^2 = (2 + 6)^2 + (-1 + 9)^2 = 64 + 64 = 128$, nên $AB = \sqrt{128} = 8\sqrt{2}$.

Đáp số: a) $\boxed{A(2; -1),\ B(-6; -9)}$; b) $\boxed{AB = 8\sqrt{2}}$
\end{solution}
\end{problem}

% ---------------------------------------------------------------- dùng biệt thức để xét số điểm chung
\begin{problem}[9DS-HSB2-038]
% nguon: dummy (tự soạn để mở rộng ngân hàng)
% lop: 9
% kc_du_kien: 9DS-HSB2-04, 9DS-HSB2-03, 9DS-PT2-02
% do_kho_du_kien: 2
Cho parabol $(P)\colon y = x^2$. Với mỗi đường thẳng sau, hãy cho biết nó cắt $(P)$ tại hai điểm, tiếp xúc với $(P)$ hay không có điểm chung với $(P)$; nếu có điểm chung thì tìm toạ độ các điểm chung đó.

a) $(d_1)\colon y = 2x + 3$;

b) $(d_2)\colon y = 4x - 4$;

c) $(d_3)\colon y = x - 1$.

\begin{solution}
Số điểm chung của đường thẳng và $(P)$ bằng số nghiệm của phương trình hoành độ giao điểm: hai nghiệm phân biệt khi $\Delta > 0$ (hay $\Delta' > 0$), nghiệm kép khi $\Delta = 0$ (đường thẳng tiếp xúc với $(P)$), vô nghiệm khi $\Delta < 0$.

a) $x^2 = 2x + 3 \Leftrightarrow x^2 - 2x - 3 = 0$, có $\Delta' = (-1)^2 - 1 \cdot (-3) = 4 > 0$, nên $(d_1)$ cắt $(P)$ tại hai điểm phân biệt.

a) Nghiệm: $x = 1 + 2 = 3$ hoặc $x = 1 - 2 = -1$; tung độ tính từ $y = x^2$ là $9$ và $1$. Hai giao điểm là $(3; 9)$ và $(-1; 1)$.

b) $x^2 = 4x - 4 \Leftrightarrow x^2 - 4x + 4 = 0$, có $\Delta' = (-2)^2 - 1 \cdot 4 = 0$, nên phương trình có nghiệm kép $x = 2$ và $(d_2)$ tiếp xúc với $(P)$ tại điểm $(2; 4)$.

c) $x^2 = x - 1 \Leftrightarrow x^2 - x + 1 = 0$, có $\Delta = (-1)^2 - 4 \cdot 1 \cdot 1 = -3 < 0$, nên phương trình vô nghiệm và $(d_3)$ không có điểm chung với $(P)$.

Đáp số: a) $\boxed{\text{Cắt tại } (3; 9) \text{ và } (-1; 1)}$; b) $\boxed{\text{Tiếp xúc tại } (2; 4)}$; c) $\boxed{\text{Không có điểm chung}}$
\end{solution}
\end{problem}

% ---------------------------------------------------------------- chứng minh đẳng thức với hàm số y = ax^2
\begin{problem}[9DS-HSB2-039]
% nguon: dummy (tự soạn để mở rộng ngân hàng)
% lop: 9
% kc_du_kien: 9DS-HSB2-01, 8DS-HDT-01
% do_kho_du_kien: 3
Cho hàm số $f(x) = \dfrac{1}{2}x^2$.

a) Chứng minh rằng $f(a + b) - f(a) - f(b) = ab$ với mọi số thực $a$, $b$.

b) Áp dụng câu a), tính $f(7) - f(3) - f(4)$ và kiểm tra lại bằng cách tính trực tiếp từng giá trị.

\begin{solution}
a) Theo công thức của hàm số: $f(a + b) = \dfrac{1}{2}(a + b)^2$, $f(a) = \dfrac{1}{2}a^2$, $f(b) = \dfrac{1}{2}b^2$.

a) Khai triển bằng hằng đẳng thức bình phương của một tổng: $f(a + b) = \dfrac{1}{2}\left(a^2 + 2ab + b^2\right)$.

a) Do đó $f(a + b) - f(a) - f(b) = \dfrac{1}{2}\left(a^2 + 2ab + b^2\right) - \dfrac{1}{2}a^2 - \dfrac{1}{2}b^2 = ab$, đúng với mọi $a$, $b$.

b) Vì $7 = 3 + 4$ nên áp dụng câu a) với $a = 3$, $b = 4$: $f(7) - f(3) - f(4) = 3 \cdot 4 = 12$.

b) Kiểm tra trực tiếp: $f(7) = \dfrac{49}{2}$, $f(3) = \dfrac{9}{2}$, $f(4) = 8$, nên $f(7) - f(3) - f(4) = \dfrac{49}{2} - \dfrac{9}{2} - 8 = 20 - 8 = 12$ (đúng).

Đáp số: a) $\boxed{f(a + b) - f(a) - f(b) = ab}$; b) $\boxed{12}$
\end{solution}
\end{problem}

% ---------------------------------------------------------------- tìm a để đường thẳng tiếp xúc với parabol
\begin{problem}[9DS-HSB2-040]
% nguon: dummy (tự soạn để mở rộng ngân hàng)
% lop: 9
% kc_du_kien: 9DS-HSB2-04, 9DS-HSB2-03
% do_kho_du_kien: 3
Tìm $a$ để đường thẳng $(d)\colon y = 2x + 1$ tiếp xúc với parabol $(P)\colon y = ax^2$ ($a \neq 0$). Tìm toạ độ tiếp điểm.

\begin{solution}
Phương trình hoành độ giao điểm của $(P)$ và $(d)$: $ax^2 = 2x + 1 \Leftrightarrow ax^2 - 2x - 1 = 0$, là phương trình bậc hai vì $a \neq 0$.

$(d)$ tiếp xúc với $(P)$ khi và chỉ khi phương trình trên có nghiệm kép, tức là $\Delta' = 0$.

Với $b' = -1$: $\Delta' = (-1)^2 - a \cdot (-1) = 1 + a$. Cho $\Delta' = 0$ được $a = -1$ (thoả mãn $a \neq 0$).

Với $a = -1$: $(P)\colon y = -x^2$ và phương trình hoành độ giao điểm là $-x^2 - 2x - 1 = 0 \Leftrightarrow (x + 1)^2 = 0$, có nghiệm kép $x = -1$.

Tung độ của tiếp điểm: $y = -(-1)^2 = -1$. Kiểm tra trên $(d)$: $2 \cdot (-1) + 1 = -1$ (đúng). Vậy tiếp điểm là $(-1; -1)$.

Đáp số: $\boxed{a = -1,\ \text{tiếp điểm } (-1; -1)}$
\end{solution}
\end{problem}

% ---------------------------------------------------------------- bài toán thực tế: tia nước của đài phun nước
\begin{problem}[9DS-HSB2-041]
% nguon: dummy (tự soạn để mở rộng ngân hàng)
% lop: 9
% kc_du_kien: 9DS-HSB2-01, 9DS-HSB2-02
% do_kho_du_kien: 3
Tia nước từ một vòi phun của đài phun nước có dạng một phần của parabol $(P)\colon y = ax^2$ trong hệ trục toạ độ $Oxy$, với gốc $O$ đặt tại đỉnh tia nước, trục $Ox$ nằm ngang, trục $Oy$ thẳng đứng hướng lên (đơn vị trên các trục là mét). Tia nước chạm mặt hồ tại điểm cách đường thẳng đứng qua $O$ là $4$ m và thấp hơn $O$ là $4$ m.

a) Tìm $a$.

b) Tại vị trí cách đường thẳng đứng qua $O$ là $2$ m, tia nước thấp hơn $O$ bao nhiêu mét?

c) Xét phần tia nước phía bên phải trục $Oy$. Tại vị trí nào của tia nước (cách đường thẳng đứng qua $O$ bao nhiêu mét) thì nó thấp hơn $O$ là $2{,}25$ m?

\begin{solution}
a) Điểm tia nước chạm mặt hồ có toạ độ $T(4; -4)$ (hoành độ $4$ vì cách trục $Oy$ là $4$ m, tung độ $-4$ vì thấp hơn $O$ là $4$ m). $T$ thuộc $(P)$ nên $-4 = a \cdot 4^2 = 16a$, suy ra $a = -\dfrac{1}{4}$.

a) Vậy $(P)\colon y = -\dfrac{1}{4}x^2$; vì $a < 0$ nên đỉnh $O$ là điểm cao nhất của tia nước, đúng với thực tế.

b) Với $x = 2$: $y = -\dfrac{1}{4} \cdot 2^2 = -1$. Vậy tại vị trí cách đường thẳng đứng qua $O$ là $2$ m, tia nước thấp hơn $O$ là $1$ m.

c) Thấp hơn $O$ là $2{,}25 = \dfrac{9}{4}$ m nghĩa là $y = -\dfrac{9}{4}$; khi đó $-\dfrac{1}{4}x^2 = -\dfrac{9}{4}$, tức là $x^2 = 9$.

c) Xét phần phía bên phải trục $Oy$ nên $x > 0$, suy ra $x = 3$; giá trị này nằm trong khoảng từ $0$ đến $4$ nên thuộc tia nước. Vậy vị trí cần tìm cách đường thẳng đứng qua $O$ là $3$ m.

Đáp số: a) $\boxed{a = -\dfrac{1}{4}}$; b) $\boxed{1 \text{ m}}$; c) $\boxed{3 \text{ m}}$
\end{solution}
\end{problem}

% ---------------------------------------------------------------- biện luận số giao điểm theo tham số
\begin{problem}[9DS-HSB2-042]
% nguon: dummy (tự soạn để mở rộng ngân hàng)
% lop: 9
% kc_du_kien: 9DS-HSB2-04, 9DS-HSB2-03
% do_kho_du_kien: 3
Cho parabol $(P)\colon y = 2x^2$ và đường thẳng $(d)\colon y = 4x + m$, với $m$ là tham số.

a) Tìm $m$ để $(d)$ cắt $(P)$ tại hai điểm phân biệt; để $(d)$ tiếp xúc với $(P)$; để $(d)$ và $(P)$ không có điểm chung.

b) Với $m = 6$, tìm toạ độ các giao điểm của $(d)$ và $(P)$.

\begin{solution}
a) Phương trình hoành độ giao điểm: $2x^2 = 4x + m \Leftrightarrow 2x^2 - 4x - m = 0$, có $\Delta' = (-2)^2 - 2 \cdot (-m) = 4 + 2m$.

a) $(d)$ cắt $(P)$ tại hai điểm phân biệt khi và chỉ khi $\Delta' > 0 \Leftrightarrow 4 + 2m > 0 \Leftrightarrow m > -2$.

a) $(d)$ tiếp xúc với $(P)$ khi và chỉ khi $\Delta' = 0 \Leftrightarrow m = -2$; $(d)$ và $(P)$ không có điểm chung khi và chỉ khi $\Delta' < 0 \Leftrightarrow m < -2$.

b) Với $m = 6$: $2x^2 - 4x - 6 = 0 \Leftrightarrow x^2 - 2x - 3 = 0$, có $\Delta' = 1 + 3 = 4$, nên $x = 1 + 2 = 3$ hoặc $x = 1 - 2 = -1$.

b) Tung độ tính từ $y = 2x^2$: với $x = 3$ thì $y = 18$; với $x = -1$ thì $y = 2$. Kiểm tra trên $(d)\colon y = 4x + 6$: $4 \cdot 3 + 6 = 18$ và $4 \cdot (-1) + 6 = 2$ (đúng).

Đáp số: a) $\boxed{m > -2 \text{ (hai điểm); } m = -2 \text{ (tiếp xúc); } m < -2 \text{ (không có điểm chung)}}$; b) $\boxed{(3; 18) \text{ và } (-1; 2)}$
\end{solution}
\end{problem}

% ---------------------------------------------------------------- hình thang tạo bởi hai giao điểm và hình chiếu trên Ox
\begin{problem}[9DS-HSB2-043]
% nguon: dummy (tự soạn để mở rộng ngân hàng)
% lop: 9
% kc_du_kien: 9DS-HSB2-03, 9DS-PT2-02
% do_kho_du_kien: 3
Cho parabol $(P)\colon y = x^2$ và đường thẳng $(d)\colon y = 2x + 8$.

a) Tìm toạ độ các giao điểm $A$, $B$ của $(d)$ và $(P)$ (hoành độ của $A$ nhỏ hơn hoành độ của $B$).

b) Gọi $H$, $K$ lần lượt là hình chiếu vuông góc của $A$, $B$ trên trục $Ox$. Tính diện tích tứ giác $ABKH$ (đơn vị đo trên các trục là xentimét).

\begin{solution}
a) Hoành độ giao điểm là nghiệm của phương trình $x^2 = 2x + 8 \Leftrightarrow x^2 - 2x - 8 = 0$.

a) Dùng công thức nghiệm thu gọn với $b' = -1$: $\Delta' = (-1)^2 - 1 \cdot (-8) = 9$, $\sqrt{\Delta'} = 3$, nên $x_1 = 1 - 3 = -2$ và $x_2 = 1 + 3 = 4$.

a) Tính tung độ từ $y = x^2$: với $x = -2$ thì $y = 4$; với $x = 4$ thì $y = 16$. Vậy $A(-2; 4)$ và $B(4; 16)$ (thử lại bằng $y = 2x + 8$: $2 \cdot (-2) + 8 = 4$, $2 \cdot 4 + 8 = 16$, đúng).

b) Hình chiếu của $A$, $B$ trên $Ox$ là $H(-2; 0)$ và $K(4; 0)$. Ta có $AH = 4$ cm, $BK = 16$ cm và $HK = 4 - (-2) = 6$ cm.

b) $AH$ và $BK$ cùng vuông góc với $Ox$ nên $AH \parallel BK$, do đó $ABKH$ là hình thang vuông với hai đáy $AH$, $BK$ và đường cao $HK$.

b) Diện tích hình thang: $S = \dfrac{(AH + BK) \cdot HK}{2} = \dfrac{(4 + 16) \cdot 6}{2} = 60$ (cm$^2$).

Đáp số: a) $\boxed{A(-2; 4),\ B(4; 16)}$; b) $\boxed{S_{ABKH} = 60 \text{ cm}^2}$
\end{solution}
\end{problem}

% ---------------------------------------------------------------- tìm a và b từ hoành độ hai giao điểm
\begin{problem}[9DS-HSB2-044]
% nguon: dummy (tự soạn để mở rộng ngân hàng)
% lop: 9
% kc_du_kien: 9DS-HSB2-03, 9DS-HSB2-01, 9DS-HPT-08
% do_kho_du_kien: 3
Biết parabol $(P)\colon y = ax^2$ và đường thẳng $(d)\colon y = 2x + b$ cắt nhau tại hai điểm có hoành độ lần lượt là $1$ và $-3$. Tìm $a$ và $b$.

\begin{solution}
Giao điểm có hoành độ $x_0$ thuộc cả $(P)$ và $(d)$, nên tung độ của nó vừa bằng $ax_0^2$ vừa bằng $2x_0 + b$, tức là $ax_0^2 = 2x_0 + b$.

Với $x_0 = 1$: $a \cdot 1^2 = 2 \cdot 1 + b$, tức là $a = 2 + b$ $(1)$.

Với $x_0 = -3$: $a \cdot (-3)^2 = 2 \cdot (-3) + b$, tức là $9a = b - 6$ $(2)$.

Từ $(1)$ suy ra $b = a - 2$; thay vào $(2)$ được $9a = a - 2 - 6$, tức là $8a = -8$, suy ra $a = -1$ và $b = -1 - 2 = -3$.

Thử lại: với $a = -1$, $b = -3$, phương trình hoành độ giao điểm $-x^2 = 2x - 3 \Leftrightarrow x^2 + 2x - 3 = 0 \Leftrightarrow (x - 1)(x + 3) = 0$ có hai nghiệm $x = 1$ và $x = -3$ (đúng), và $a = -1 \neq 0$.

Đáp số: $\boxed{a = -1,\ b = -3}$
\end{solution}
\end{problem}

% ---------------------------------------------------------------- tham số, định lý Viète, tổng bình phương hai hoành độ
\begin{problem}[9DS-HSB2-045]
% nguon: dummy (tự soạn để mở rộng ngân hàng)
% lop: 9
% kc_du_kien: 9DS-HSB2-04, 9DS-HSB2-03, 9DS-PT2-03
% do_kho_du_kien: 4
Cho parabol $(P)\colon y = 2x^2$ và đường thẳng $(d)\colon y = 2x + m$, với $m$ là tham số.

a) Tìm $m$ để $(d)$ cắt $(P)$ tại hai điểm phân biệt.

b) Gọi $x_1$, $x_2$ là hoành độ hai giao điểm đó. Tìm $m$ để $x_1^2 + x_2^2 = 5$.

c) Với giá trị $m$ tìm được ở câu b), tìm toạ độ các giao điểm.

\begin{solution}
a) Phương trình hoành độ giao điểm: $2x^2 = 2x + m \Leftrightarrow 2x^2 - 2x - m = 0$, có $\Delta' = (-1)^2 - 2 \cdot (-m) = 1 + 2m$.

a) $(d)$ cắt $(P)$ tại hai điểm phân biệt khi và chỉ khi $\Delta' > 0 \Leftrightarrow 1 + 2m > 0 \Leftrightarrow m > -\dfrac{1}{2}$.

b) Với $m > -\dfrac{1}{2}$, theo định lý Viète: $x_1 + x_2 = \dfrac{2}{2} = 1$ và $x_1x_2 = \dfrac{-m}{2}$.

b) Biến đổi: $x_1^2 + x_2^2 = (x_1 + x_2)^2 - 2x_1x_2 = 1^2 - 2 \cdot \dfrac{-m}{2} = 1 + m$.

b) Cho $x_1^2 + x_2^2 = 5$: $1 + m = 5$, suy ra $m = 4$ (thoả mãn $m > -\dfrac{1}{2}$).

c) Với $m = 4$: $2x^2 - 2x - 4 = 0 \Leftrightarrow x^2 - x - 2 = 0$, có $a - b + c = 1 + 1 - 2 = 0$ nên $x = -1$ hoặc $x = -\dfrac{c}{a} = 2$.

c) Tung độ tính từ $y = 2x^2$: với $x = -1$ thì $y = 2$; với $x = 2$ thì $y = 8$. Vậy hai giao điểm là $(-1; 2)$ và $(2; 8)$. Kiểm tra: $(-1)^2 + 2^2 = 5$ (đúng).

Đáp số: a) $\boxed{m > -\dfrac{1}{2}}$; b) $\boxed{m = 4}$; c) $\boxed{(-1; 2) \text{ và } (2; 8)}$
\end{solution}
\end{problem}

% ---------------------------------------------------------------- hai điểm đối xứng qua Oy tạo với O tam giác đều
\begin{problem}[9DS-HSB2-046]
% nguon: dummy (tự soạn để mở rộng ngân hàng)
% lop: 9
% kc_du_kien: 9DS-HSB2-02, 9DS-HSB2-01
% do_kho_du_kien: 4
Cho parabol $(P)\colon y = x^2$. Tìm toạ độ hai điểm $A$, $B$ thuộc $(P)$, đối xứng với nhau qua trục $Oy$, sao cho tam giác $OAB$ là tam giác đều. Tính diện tích tam giác $OAB$.

\begin{solution}
Parabol $(P)$ nhận trục $Oy$ làm trục đối xứng. Gọi $B(t; t^2)$ với $t > 0$ thì điểm đối xứng với $B$ qua $Oy$ là $A(-t; t^2)$.

Hai điểm $A$, $B$ có cùng tung độ nên $AB = t - (-t) = 2t$. Ngoài ra $OA = OB = \sqrt{t^2 + t^4}$ (hai điểm đối xứng qua $Oy$ cách đều $O$), nên tam giác $OAB$ cân tại $O$.

Tam giác cân $OAB$ là tam giác đều khi và chỉ khi $OA = AB$, tức là $t^2 + t^4 = 4t^2$, hay $t^4 = 3t^2$.

Vì $t > 0$ nên chia hai vế cho $t^2$ được $t^2 = 3$, suy ra $t = \sqrt{3}$. Vậy $A\left(-\sqrt{3}; 3\right)$ và $B\left(\sqrt{3}; 3\right)$. Kiểm tra: $OA^2 = 3 + 9 = 12 = AB^2$ (đúng).

Cạnh $AB = 2\sqrt{3}$; đường cao kẻ từ $O$ xuống $AB$ (nằm trên đường thẳng $y = 3$) bằng $3$. Diện tích: $S = \dfrac{1}{2} \cdot AB \cdot 3 = \dfrac{1}{2} \cdot 2\sqrt{3} \cdot 3 = 3\sqrt{3}$.

Đáp số: $\boxed{A\left(-\sqrt{3}; 3\right),\ B\left(\sqrt{3}; 3\right),\ S_{OAB} = 3\sqrt{3}}$
\end{solution}
\end{problem}

% ---------------------------------------------------------------- hai giao điểm khác phía trục tung, tỉ số hai hoành độ, diện tích
\begin{problem}[9DS-HSB2-047]
% nguon: dummy (tự soạn để mở rộng ngân hàng)
% lop: 9
% kc_du_kien: 9DS-HSB2-04, 9DS-HSB2-03, 9DS-PT2-03
% do_kho_du_kien: 4
Cho parabol $(P)\colon y = \dfrac{1}{3}x^2$ và đường thẳng $(d)\colon y = mx + 6$, với $m$ là tham số.

a) Chứng minh rằng với mọi $m$, $(d)$ luôn cắt $(P)$ tại hai điểm phân biệt $A$, $B$ nằm về hai phía của trục tung.

b) Gọi $x_1$, $x_2$ là hoành độ của $A$, $B$. Tìm $m$ để $|x_1| = 2|x_2|$.

c) Với $m = 1$, tính diện tích tam giác $OAB$ ($O$ là gốc toạ độ, đơn vị đo trên các trục là xentimét).

\begin{solution}
a) Phương trình hoành độ giao điểm: $\dfrac{1}{3}x^2 = mx + 6$; nhân hai vế với $3$ rồi chuyển vế: $x^2 - 3mx - 18 = 0$.

a) Có $\Delta = (-3m)^2 - 4 \cdot 1 \cdot (-18) = 9m^2 + 72 > 0$ với mọi $m$, nên phương trình luôn có hai nghiệm phân biệt, tức là $(d)$ luôn cắt $(P)$ tại hai điểm phân biệt.

a) Theo định lý Viète: $x_1 + x_2 = 3m$ và $x_1x_2 = -18 < 0$, nên $x_1$, $x_2$ trái dấu; vậy $A$, $B$ nằm về hai phía của trục tung.

b) Vì $x_1$, $x_2$ trái dấu nên $|x_1| = 2|x_2|$ khi và chỉ khi $x_1 = -2x_2$. Thay vào $x_1x_2 = -18$ được $-2x_2^2 = -18$, tức là $x_2^2 = 9$, suy ra $x_2 = 3$ hoặc $x_2 = -3$.

b) Nếu $x_2 = 3$ thì $x_1 = -6$ và $x_1 + x_2 = -3 = 3m$, suy ra $m = -1$. Nếu $x_2 = -3$ thì $x_1 = 6$ và $x_1 + x_2 = 3 = 3m$, suy ra $m = 1$.

b) Thử lại: $m = 1$ cho $x^2 - 3x - 18 = (x - 6)(x + 3) = 0$, hai nghiệm $6$ và $-3$ thoả $|6| = 2 \cdot |-3|$; $m = -1$ cho $x^2 + 3x - 18 = (x + 6)(x - 3) = 0$, hai nghiệm $-6$ và $3$ thoả $|-6| = 2 \cdot |3|$.

c) Với $m = 1$, hai nghiệm là $x = -3$ và $x = 6$; tung độ tính từ $y = \dfrac{1}{3}x^2$ là $3$ và $12$. Vậy $A(-3; 3)$ và $B(6; 12)$ (thử lại bằng $y = x + 6$: $-3 + 6 = 3$, $6 + 6 = 12$, đúng).

c) $(d)\colon y = x + 6$ cắt trục tung tại $C(0; 6)$, nên $OC = 6$ cm. Vì $A$, $B$ nằm về hai phía của trục tung nên $C$ nằm giữa $A$ và $B$, do đó $S_{OAB} = S_{OAC} + S_{OBC}$.

c) Khoảng cách từ $A$, $B$ đến trục tung là $3$ cm và $6$ cm, nên $S_{OAC} = \dfrac{1}{2} \cdot 6 \cdot 3 = 9$ (cm$^2$) và $S_{OBC} = \dfrac{1}{2} \cdot 6 \cdot 6 = 18$ (cm$^2$). Suy ra $S_{OAB} = 9 + 18 = 27$ (cm$^2$).

Đáp số: a) $\boxed{\Delta = 9m^2 + 72 > 0 \text{ và } x_1x_2 = -18 < 0 \text{ với mọi } m}$; b) $\boxed{m = 1 \text{ hoặc } m = -1}$; c) $\boxed{S_{OAB} = 27 \text{ cm}^2}$
\end{solution}
\end{problem}

% ---------------------------------------------------------------- bài toán thực tế: dây cáp cầu treo
\begin{problem}[9DS-HSB2-048]
% nguon: dummy (tự soạn để mở rộng ngân hàng)
% lop: 9
% kc_du_kien: 9DS-HSB2-01, 9DS-HSB2-02
% do_kho_du_kien: 4
Dây cáp của một cây cầu treo có dạng một phần của parabol. Hai đầu dây cáp buộc vào đỉnh của hai tháp cầu; mỗi tháp cao $20$ m tính từ mặt cầu và hai tháp cách nhau $120$ m. Điểm thấp nhất của dây cáp nằm chính giữa hai tháp và cách mặt cầu $2$ m. Chọn hệ trục toạ độ $Oxy$ có gốc $O$ tại điểm thấp nhất của dây cáp, trục $Ox$ song song với mặt cầu, trục $Oy$ thẳng đứng hướng lên (đơn vị trên các trục là mét); khi đó dây cáp là một phần của parabol $(P)\colon y = ax^2$.

a) Tìm $a$.

b) Các thanh treo thẳng đứng nối dây cáp xuống mặt cầu. Tính chiều dài thanh treo đặt cách điểm chính giữa $20$ m.

c) Một thanh treo khác dài $10$ m. Thanh treo đó cách điểm chính giữa bao nhiêu mét?

\begin{solution}
a) Điểm thấp nhất của dây cáp là gốc $O$ và cách mặt cầu $2$ m, nên mặt cầu nằm trên đường thẳng $y = -2$. Đỉnh mỗi tháp cao $20$ m so với mặt cầu nên có tung độ $20 - 2 = 18$.

a) Hai tháp cách nhau $120$ m và đối xứng qua trục $Oy$ nên hoành độ chân tháp là $\pm 60$. Đỉnh tháp $(60; 18)$ thuộc $(P)$: $18 = a \cdot 60^2 = 3600a$, suy ra $a = \dfrac{1}{200}$.

b) Tại vị trí cách điểm chính giữa $20$ m, hoành độ là $x = 20$ (hoặc $-20$) và điểm của dây cáp có tung độ $y = \dfrac{1}{200} \cdot 20^2 = 2$.

b) Mặt cầu nằm trên đường thẳng $y = -2$ nên chiều dài thanh treo bằng hiệu hai tung độ: $2 - (-2) = 4$ (m).

c) Thanh treo dài $10$ m nên điểm của dây cáp có tung độ $y$ thoả mãn $y - (-2) = 10$, suy ra $y = 8$.

c) Cho $\dfrac{1}{200}x^2 = 8$ được $x^2 = 1\,600$, suy ra $x = 40$ hoặc $x = -40$. Vậy thanh treo dài $10$ m cách điểm chính giữa $40$ m (về mỗi phía).

Đáp số: a) $\boxed{a = \dfrac{1}{200}}$; b) $\boxed{4 \text{ m}}$; c) $\boxed{40 \text{ m}}$
\end{solution}
\end{problem}

% ---------------------------------------------------------------- điểm trên parabol gần một điểm cố định nhất
\begin{problem}[9DS-HSB2-049]
% nguon: dummy (tự soạn để mở rộng ngân hàng)
% lop: 9
% kc_du_kien: 9DS-HSB2-01, 9DS-HSB2-02, 8DS-HDT-01
% do_kho_du_kien: 4
Cho parabol $(P)\colon y = x^2$ và điểm $A\left(0; \dfrac{3}{2}\right)$. Điểm $M$ chạy trên $(P)$. Tìm toạ độ điểm $M$ để độ dài $AM$ nhỏ nhất và tính giá trị nhỏ nhất đó.

\begin{solution}
Gọi $M(x; y)$ thuộc $(P)$, khi đó $y = x^2$ và $y \geq 0$. Theo công thức khoảng cách giữa hai điểm: $AM^2 = x^2 + \left(y - \dfrac{3}{2}\right)^2$.

Thay $x^2 = y$: $AM^2 = y + y^2 - 3y + \dfrac{9}{4} = y^2 - 2y + \dfrac{9}{4}$.

Dùng hằng đẳng thức bình phương của một hiệu: $AM^2 = (y - 1)^2 + \dfrac{5}{4}$.

Vì $(y - 1)^2 \geq 0$ nên $AM^2 \geq \dfrac{5}{4}$, suy ra $AM \geq \dfrac{\sqrt{5}}{2}$. Dấu bằng xảy ra khi và chỉ khi $y = 1$ (thoả mãn $y \geq 0$).

Khi $y = 1$ thì $x^2 = 1$, suy ra $x = 1$ hoặc $x = -1$. Thử lại với $M(1; 1)$: $AM^2 = 1^2 + \left(1 - \dfrac{3}{2}\right)^2 = 1 + \dfrac{1}{4} = \dfrac{5}{4}$ (đúng); điểm $M(-1; 1)$ đối xứng với $M(1; 1)$ qua $Oy$ nên cho cùng kết quả.

Đáp số: $\boxed{M(1; 1) \text{ hoặc } M(-1; 1),\ AM_{\min} = \dfrac{\sqrt{5}}{2}}$
\end{solution}
\end{problem}

% ---------------------------------------------------------------- hai điểm của parabol, tam giác vuông tại gốc toạ độ, điểm cố định
\begin{problem}[9DS-HSB2-050]
% nguon: dummy (tự soạn để mở rộng ngân hàng)
% lop: 9
% kc_du_kien: 9DS-HSB2-01, 9DS-HSB2-02, 9DS-HSB2-03
% do_kho_du_kien: 5
Cho parabol $(P)\colon y = x^2$. Hai điểm phân biệt $A$, $B$ khác gốc toạ độ $O$ thuộc $(P)$ có hoành độ lần lượt là $a$ và $b$.

a) Chứng minh rằng tam giác $OAB$ vuông tại $O$ khi và chỉ khi $ab = -1$.

b) Cho $A$ có hoành độ bằng $2$. Tìm toạ độ điểm $B$ để tam giác $OAB$ vuông tại $O$ và tính diện tích tam giác $OAB$.

c) Chứng minh rằng khi tam giác $OAB$ vuông tại $O$ thì đường thẳng $AB$ luôn đi qua một điểm cố định.

\begin{solution}
a) Vì $A$, $B$ thuộc $(P)$ nên $A(a; a^2)$ và $B(b; b^2)$ với $a \neq 0$, $b \neq 0$, $a \neq b$. Ta có $OA^2 = a^2 + a^4$ và $OB^2 = b^2 + b^4$.

a) Tính $AB^2 = (a - b)^2 + \left(a^2 - b^2\right)^2 = a^2 - 2ab + b^2 + a^4 - 2a^2b^2 + b^4$.

a) Theo định lí Pythagore và định lí đảo, tam giác $OAB$ vuông tại $O$ khi và chỉ khi $OA^2 + OB^2 = AB^2$, tức là $a^2 + a^4 + b^2 + b^4 = a^2 - 2ab + b^2 + a^4 - 2a^2b^2 + b^4$.

a) Rút gọn hai vế được $0 = -2ab - 2a^2b^2$, tức là $2ab(1 + ab) = 0$. Vì $a \neq 0$, $b \neq 0$ nên $ab \neq 0$, do đó điều kiện tương đương với $1 + ab = 0$, tức là $ab = -1$.

b) Với $a = 2$: điều kiện $ab = -1$ cho $b = -\dfrac{1}{2}$, nên $A(2; 4)$ và $B\left(-\dfrac{1}{2}; \dfrac{1}{4}\right)$.

b) Tính các cạnh góc vuông: $OA^2 = 4 + 16 = 20$, nên $OA = 2\sqrt{5}$; $OB^2 = \dfrac{1}{4} + \dfrac{1}{16} = \dfrac{5}{16}$, nên $OB = \dfrac{\sqrt{5}}{4}$.

b) Kiểm tra: $AB^2 = \left(2 + \dfrac{1}{2}\right)^2 + \left(4 - \dfrac{1}{4}\right)^2 = \dfrac{25}{4} + \dfrac{225}{16} = \dfrac{325}{16}$ và $OA^2 + OB^2 = 20 + \dfrac{5}{16} = \dfrac{325}{16}$ (đúng).

b) Diện tích tam giác vuông bằng nửa tích hai cạnh góc vuông: $S_{OAB} = \dfrac{1}{2} \cdot OA \cdot OB = \dfrac{1}{2} \cdot 2\sqrt{5} \cdot \dfrac{\sqrt{5}}{4} = \dfrac{5}{4}$.

c) Đường thẳng $AB$ có dạng $y = kx + n$, đi qua $A(a; a^2)$ và $B(b; b^2)$ nên $a^2 = ka + n$ và $b^2 = kb + n$. Trừ vế theo vế: $a^2 - b^2 = k(a - b)$, suy ra $k = a + b$ (vì $a \neq b$).

c) Từ $a^2 = (a + b)a + n$ suy ra $n = -ab$. Khi $OAB$ vuông tại $O$ thì $ab = -1$ nên $n = 1$ và $AB\colon y = (a + b)x + 1$.

c) Đường thẳng $y = (a + b)x + 1$ đi qua điểm $(0; 1)$ với mọi giá trị của $a + b$, nên $AB$ luôn đi qua điểm cố định $(0; 1)$.

Đáp số: a) $\boxed{OAB \text{ vuông tại } O \Leftrightarrow ab = -1}$; b) $\boxed{B\left(-\dfrac{1}{2}; \dfrac{1}{4}\right),\ S_{OAB} = \dfrac{5}{4}}$; c) $\boxed{\text{Điểm cố định } (0; 1)}$
\end{solution}
\end{problem}

% ---------------------------------------------------------------- tìm a để tam giác tạo bởi gốc toạ độ và hai giao điểm có diện tích cho trước
\begin{problem}[9DS-HSB2-051]
% nguon: dummy (tự soạn để mở rộng ngân hàng)
% lop: 9
% kc_du_kien: 9DS-HSB2-04, 9DS-HSB2-03, 9DS-PT2-03
% do_kho_du_kien: 5
Cho parabol $(P)\colon y = ax^2$ ($a \neq 0$) và đường thẳng $(d)\colon y = x + 2$. Biết $(d)$ cắt $(P)$ tại hai điểm phân biệt $A$, $B$. Tìm $a$ để tam giác $OAB$ có diện tích bằng $3$ ($O$ là gốc toạ độ, đơn vị diện tích).

\begin{solution}
Phương trình hoành độ giao điểm: $ax^2 = x + 2 \Leftrightarrow ax^2 - x - 2 = 0$, có $\Delta = 1 + 8a$. $(d)$ cắt $(P)$ tại hai điểm phân biệt khi và chỉ khi $\Delta > 0$, tức là $a > -\dfrac{1}{8}$ (và $a \neq 0$).

Gọi $x_1$, $x_2$ là hoành độ của $A$, $B$. Theo định lý Viète: $x_1 + x_2 = \dfrac{1}{a}$ và $x_1x_2 = -\dfrac{2}{a}$.

$(d)$ cắt trục tung tại $C(0; 2)$, nên $OC = 2$; khoảng cách từ $A$, $B$ đến trục tung lần lượt là $|x_1|$ và $|x_2|$.

Nếu $A$, $B$ nằm khác phía của trục tung (khi $x_1$, $x_2$ trái dấu) thì $S_{OAB} = S_{OAC} + S_{OBC} = \dfrac{1}{2} \cdot 2 \cdot |x_1| + \dfrac{1}{2} \cdot 2 \cdot |x_2| = |x_1| + |x_2| = |x_1 - x_2|$.

Nếu $A$, $B$ nằm cùng phía của trục tung (khi $x_1$, $x_2$ cùng dấu) thì $S_{OAB} = \left|S_{OAC} - S_{OBC}\right| = \left||x_1| - |x_2|\right| = |x_1 - x_2|$. Vậy trong cả hai trường hợp $S_{OAB} = |x_1 - x_2|$.

Tính $(x_1 - x_2)^2 = (x_1 + x_2)^2 - 4x_1x_2 = \dfrac{1}{a^2} + \dfrac{8}{a} = \dfrac{1 + 8a}{a^2}$.

Điều kiện $S_{OAB} = 3$ cho $\dfrac{1 + 8a}{a^2} = 9$, tức là $9a^2 - 8a - 1 = 0$. Phân tích: $(a - 1)(9a + 1) = 0$, nên $a = 1$ hoặc $a = -\dfrac{1}{9}$.

Đối chiếu điều kiện $a > -\dfrac{1}{8}$, $a \neq 0$: $1 > -\dfrac{1}{8}$ và $-\dfrac{1}{9} > -\dfrac{1}{8}$, nên cả hai giá trị đều thoả mãn.

Thử lại: với $a = 1$, $x^2 - x - 2 = 0$ có nghiệm $-1$ và $2$, $A(-1; 1)$, $B(2; 4)$ khác phía và $S = \dfrac{1}{2} \cdot 2 \cdot (1 + 2) = 3$; với $a = -\dfrac{1}{9}$, $x^2 + 9x + 18 = 0$ có nghiệm $-3$ và $-6$, $A(-3; -1)$, $B(-6; -4)$ cùng phía và $S = \dfrac{1}{2} \cdot 2 \cdot (6 - 3) = 3$.

Đáp số: $\boxed{a = 1 \text{ hoặc } a = -\dfrac{1}{9}}$
\end{solution}
\end{problem}

\end{document}
